% 前置部分：内容简介 + 图书在版编目（CIP）数据 / 版权页 + 第一版前言

% ---------- 内容简介 + CIP 版权页（不进目录） ----------
\begingroup
\pagestyle{empty}

{\headfont\zihao{-4} 内容简介\par}
\vspace{0.5em}

本书是教育部``高等教育面向 $21$ 世纪教学内容和课程体系改革计划''的研究成果，是面向 $21$ 世纪课程教材。教程用``连续量的演算体系及其数学理论''的全新观点统率全书，在保留传统数学分析基本内容的前提下，比较好地处理极限与微积分演算及应用的关系，建立了一个既循序渐进、生动直观，又保持了严密性的系统，与传统的教程十分不同。本教程对概念、方法的来源与实质，有许多独到的、精辟的见解。教程分上、下两册，本书为上册，主要内容包括实数连续统、函数、极限与函数连续性、微商与微分、微分中值定理及其应用、不定积分、定积分、微积分进一步应用、再论实数系等。本书是作者集几十年教学与教改经验之力作，在教学改革实践中取得较好的效果。

本书可作为高等学校理科及师范学校数学学科各专业的教科书，也可供计算机学科、力学、物理学科各专业选用及社会读者阅读。

\vspace{1em}
{\headfont\zihao{-4} 图书在版编目（CIP）数据\par}
\vspace{0.3em}

数学分析简明教程．上册／邓东皋，尹小玲编著．— $2$ 版．— 北京：高等教育出版社，$2006.3$

ISBN $7-04-018662-4$

{\rmfamily I}．数… {\rmfamily II}．\textcircled{\raisebox{-0.15ex}{\scriptsize 1}}邓… \textcircled{\raisebox{-0.15ex}{\scriptsize 2}}尹… {\rmfamily III}．数学分析－高等学校－教材 {\rmfamily IV}．$017$

中国版本图书馆 CIP 数据核字（$2006$）第 $006000$ 号

\vspace{0.8em}
\noindent
\begin{minipage}[t]{0.32\textwidth}
\zihao{-5}
策划编辑\quad 李艳馥\\
责任编辑\quad 黄建英\\
责任编辑\quad 张长虹
\end{minipage}%
\begin{minipage}[t]{0.32\textwidth}
\zihao{-5}
封面设计\quad 张\hspace{1em}楠\\
版式设计\quad 范晓红
\end{minipage}%
\begin{minipage}[t]{0.32\textwidth}
\zihao{-5}
责任印制\quad 杨\hspace{1em}明\\
责任校对\quad 金\hspace{1em}辉
\end{minipage}

\vspace{0.8em}
{\color{black}\hrule height0.8pt}
\vspace{0.8em}

\noindent
\begin{minipage}[t]{0.5\textwidth}
\zihao{-5}
\begin{tabular}{@{}ll@{}}
出版发行 & 高等教育出版社\\
社\quad\quad 址 & 北京市西城区德外大街\,$4$\,号\\
邮政编码 & $100011$\\
总\quad\quad 机 & $010-58581000$\\[0.6em]
经\quad\quad 销 & 蓝色畅想图书发行有限公司\\
印\quad\quad 刷 & 国防工业出版社印刷厂\\[0.6em]
开\quad\quad 本 & $787\times960$\ \ $1/16$\\
印\quad\quad 张 & $20.5$\\
字\quad\quad 数 & $370\ 000$
\end{tabular}
\end{minipage}%
\begin{minipage}[t]{0.5\textwidth}
\zihao{-5}
\begin{tabular}{@{}ll@{}}
购书热线 & $010-58581118$\\
免费咨询 & $800-810-0598$\\
网\quad\quad 址 & \nolinkurl{http://www.hep.edu.cn}\\
 & \nolinkurl{http://www.hep.com.cn}\\[0.6em]
网上订购 & \nolinkurl{http://www.landraco.com}\\
 & \nolinkurl{http://www.landraco.com.cn}\\
畅想教育 & \nolinkurl{http://www.widedu.com}\\[0.6em]
版\quad\quad 次 & $1999$ 年 $6$ 月第 $1$ 版\\
 & $2006$ 年 $3$ 月第 $2$ 版\\
印\quad\quad 次 & $2006$ 年 $3$ 月第 $1$ 次印刷\\
定\quad\quad 价 & $23.70$ 元
\end{tabular}
\end{minipage}

\vspace{0.8em}
{\color{black}\hrule height0.8pt}
\vspace{0.6em}

{\zihao{-5} 本书如有缺页、倒页、脱页等质量问题，请到所购图书销售部门联系调换。\par}
{\headfont\zihao{-5} 版权所有\quad 侵权必究\par}
{\bfseries\zihao{-5} 物料号\quad $18662$-$00$\par}
\clearpage
\endgroup

% ---------- 第一版前言（进目录） ----------
\chapter*{第一版前言}
\addcontentsline{toc}{chapter}{第一版前言}

数学分析的主要内容是微积分，这是人类在科学中最伟大的创造之一。微积分研究的对象是连续量。本教程提供给读者的是一个连续量的演算体系及其数学理论。过去读者在中小学学的算术与代数的演算大都只涉及离散量，本教程将提供一套崭新的演算——连续量的演算。一个连续量对另一个连续量的连续依赖，其基本问题之一是``瞬时''变化率，或一个连续量对另一个连续量的变化``速率''，这就引导到微商的概念。变化率要``瞬时''，这是连续量的特征之一。变化率为什么要``瞬时''，其根本原因是，这样就能``机械化''地进行演算了。另一个基本问题是连续变化的积累，或连续作用的总和。这就引导到积分的概念。牛顿与莱布尼茨在创立微积分时的重大贡献之一是发现求这种连续量作用的积累或总和，是求变化率运算的逆运算，从而建立了一套连续量的``机械化''的演算体系。这一切最重要的体现是立微分方程与解微分方程。实数本质上是（一维）连续量的数学模型。本教程上册讲的一元函数微积分实际上是初等函数微积分。为了把它推广到非初等函数，人们才需要无穷级数与含参变量积分这样的工具，同时为了解决多个连续量之间的依赖关系问题，才需要发展到多元微积分。后面这两部分（无穷级数与多元微积分）便构成了本教程下册的主要内容。极限是对上述所有概念形式化统一处理的工具。用极限可以把上述概念精确化和统一处理，使理论简明统一。因此，极限的概念与运算将贯穿全书。但应提醒读者注意，一方面不要因为极限贯穿全书便用它掩盖了数学分析研究连续量演算体系的本质；另一方面，对极限的掌握也是通过对微积分各项内容的研究而逐步加深的。这是一个循序渐进的过程，读者不必希望``一蹴而就''。

历史上，微积分运算体系形成在先，用极限与实数理论给运算体系建立严格的数学基础在后。如何处理这两者的关系，一直是数学分析教材体系的中心问题。过去，说得极端一些，大体有两种讲法。一种是把整个教材分成两部分：初等微积分与高等微积分。在讲初等微积分时，主要讲微积分的运算与应用，在讲高等微积分时才讲严格的理论。另一种讲法是先把实数、极限、连续等都讲述清楚，然后再把微积分、级数等看成不同类型的极限。由于前面讲授极限的篇幅很大，因此通常把后面这种讲法称作``大头分析''。这两种体系各有优缺点，且曾经分别在不同时期在我国流行过。本教程的编写，是希望在上述``初高等微积分''与``大头分析''两种讲法之间，走出一条真正的循序渐进之路，而整个体系又是逻辑上完整的。为此，本书把实数连续性放在一个基本的位置，根据戴德金的原始思想，给出判断一个有序稠密数系是否连续的``戴德金连续性判别准则''。然后，通过实数的十进表示证明实数系是连续的。在开头讲极限时，严格地不出现上下确界、子序列、一致连续、哥西准则等等概念。闭区间连续函数的每个性质，只有在它需要时才分别出现，并用实数连续性定理直接加以证明。在一元微积分讲完以后，再综合讲实数连续性的其他描述与实数系的拓扑性质。全文下来每一步都是严格的。我们这样处理，与别的教程十分不同，其目的是使读者初学时尽量减少新的概念，并使微积分的运算体系与应用能尽快出现。

数学分析讲授的最大倾向是过于形式化。本教程的另一项努力的目标是在讲授数学内容时，把形式逻辑推导所掩盖的概念的背景来源、解决问题的思考方法、所讲授的内容在整个理论体系中的作用与地位以及和其他概念之间的内在联系等揭示出来。本书在讲述过程中，尽可能使读者注意物理与力学的背景模型（实物），几何的形象直观（形象），抽象的演算推理（数量）三者的结合，把它们融为一体。所有这些，人们有时用``揭露实质''这个词来解释它。当然，不能要求所有的数学书都这样写，但作为理科（包括师范）数学系一门主要的基础课，是大学生第一次接触高等数学的地方，我们认为这种写法是应该提倡的。

本书保留了传统数学分析的基本内容。我们认为这些基本概念、基本理论与基本训练，对理科数学系的学生是必不可少的。但我们努力使选材的确符合``基本''的要求，而不过分展开，使本书成为一本名副其实的``简明教程''。在与近代数学接轨的问题上，本书更多地体现在内容叙述的观点上，而不强调增加新内容。但和国内传统数学分析教材不同的是，本书讲述了一点常微分方程，并用微积分从开普勒三定律推导出万有引力定律，以及反过来从万有引力定律推导出开普勒三定律。我们认为，为引起读者的兴趣，体会微积分演算体系的精神实质，花不多的时间去讲授，是完全值得的。

在使用本书时，一些教学内容的先后次序是可以根据需要加以调整的。如上册的第八章与第九章，下册的级数理论与多元微积分，先讲后讲都是可以变更的。在上册教学中，只有两个定理（闭区间连续函数的最值定理与一致连续性定理）的证明稍微困难。有些高校在教学时可根据实际情况，在第一次讲授时把证明略去（因为后面第九章还有别的证明）。还有，下册中傅里叶级数的平均收敛性与重积分换元公式的证明，不同学校也可酌情取舍。

本书第一位编者曾长期在北京大学工作，讲授过多年的数学分析课程。在教学中，曾得到过北京大学许多老一辈数学家的指导，获益甚多。他们中有申又枨、闵嗣鹤、程民德、吴光磊、冷生明，编者对他们表示衷心的感谢。其中特别是吴光磊先生与闵嗣鹤先生，他们曾帮助编者提高了对许多数学问题的认识，本书中的许多看法都曾受到过他们的启发。不幸他们俩已去世，本书的出版也是对他们的纪念。本书于去年 $7$ 月作为讲义打印出来后，中山大学林伟教授、范达教授、周建伟教授、朱匀华教授、丘兆福副教授、姚正安副教授、胡建勋副教授、北京大学文丽教授、北京师范大学王昆扬教授以及陕西师范大学常心怡教授，都阅读了或试用过本书，他们对本书提出了许多宝贵的意见，我们对他们也表示衷心的感谢。另外，中山大学校领导与教务处，对本书的出版给予了很大的关心与支持，我们衷心地感谢他们。

本书成书的时间不长，疏漏错误难免，敬请读者批评指正。

\vspace{2em}
\begin{flushright}
邓东皋\quad 尹小玲\\[0.5em]
$1998$ 年 $11$ 月 $12$ 日\\[0.5em]
中山大学蒲园
\end{flushright}
\clearpage
